% 子文件：6.定积分.tex
% 主文档中使用 \input{6.定积分} 引入，不含 document 环境

\section{定积分的概念与性质}
\subsection{定积分的概念}

定义6.1\quad 设函数$f(x)$在区间$[a,b]$上有定义，用分点
\[
a=x_1<x_2<\cdots<x_i<x_{i+1}<\cdots<x_n<x_{n+1}=b
\]
将区间$[a,b]$分为$n$个小区间$[x_i,x_{i+1}]$，记$\Delta x_i=x_{i+1}-x_i$，$\lambda=\max\{|\Delta x_i|\}$。任取$\xi_i\in[x_i,x_{i+1}],\ i=1,2,\cdots,n$。如果乘积的和式
\[
\sum_{i=1}^{n}f(\xi_i)\Delta x_i
\]
（称为积分和）的极限
\[
\lim_{\lambda\to0}\sum_{i=1}^{n}f(\xi_i)\Delta x_i
\]
存在，且这个极限与$x_i$和$\xi_i$的取法无关，则称$f(x)$在$[a,b]$上可积，并称此极限值为$f(x)$在区间$[a,b]$上由$a$到$b$的定积分，用记号$\int _a^bf(x)dx$表示之，即
\[
\int _a^bf(x)dx=\lim_{\lambda\to0}\sum_{i=1}^{n}f(\xi_i)\Delta x_i.
\]

定理6.1\quad 如果$f(x)$在$[a,b]$上可积，则$f(x)$在$[a,b]$上必有界。

定理6.2\quad 如果

定理6.3\quad 如果

\subsection{定积分的简单性质}
定积分是由被积函数与积分区间所确定的一个数。由此不难得到下列性质。在本段中，假定所涉及的定积分都存在。
\begin{enumerate}
    \item 
    \item 
    \item 
    \item $\int_a^b\bigl[kf(x)+lg(x)\bigr]dx=k\int_a^bf(x)dx+l\int_a^bg(x)dx\quad(k,l为常数).$
    \item 
    \item 
    \item 
    \item 
    \item 
    \item 定积分中值定理\quad 设$f(x)\in C[a,b]$，则至少存在一点$\xi\in[a,b]$，使得

上式也叫做平均值公式。求连续变量的平均值就要用到它。
\end{enumerate}

\noindent 设，证明柯西不等式。

\section{微积分基本定理}

定理6.4

它还说明，连续函数$f(x)$

定理6.5\quad 如果$F(x)$是$[a,b]$区间上连续函数$f(x)$的一个原函数，则
\[
\int_a^bf(x)dx=F(b)-F(a) \tag{3}
\]

证明\quad 因为

$C$是待定常数，即有

令$x=a$

\noindent 特别，公式(3)称为牛顿—莱布尼茨公式。

\noindent 公式

\section{定积分的计算}
\subsection{定积分的换元积分法}

定理6.6\quad 设$f(x)\in C[a,b]$，对变换$x=\phi(t)$，若有$\alpha,\beta$满足：
\begin{enumerate}
    \item[(i)] $\phi(\alpha)=a,\ \phi(\beta)=b$；
    \item[(ii)] 在$\alpha,\beta$界定的区间上，$a\le \phi(t)\le b$；
    \item[(iii)] 在$\alpha,\beta$界定的区间上，$\phi(t)$有连续的导数，
\end{enumerate}
则
\[
\int_a^bf(x)dx=\int_{\alpha}^{\beta}f\bigl[\phi(t)\bigr]\phi'(t)dt.
\]
这个定理说明用换元积分法计算定积分时，应把积分上、下限同时换为新的积分变量的上、下限，通过新的积分算出积分值。这样避免了求$f(x)$的原函数，所以对变换$x=\phi(t)$也不要求它有反函数。

设$f(x)$在区间

由定积分定义不难推证，对一般可积函数

\subsection{定积分的分部积分法}

定理6.7\quad 设$u(x),v(x)$在区间$[a,b]$上有连续的导数，则

由不定积分的分部积分法以及牛顿—莱布尼茨公式，这是显然的。

\section{反常积分}
定积分$\int_a^bf(x)dx$受到两个限制，其一，积分区间$[a,b]$是有限区间；其二，被积函数在积分区间上是有界函数。

\subsection{无穷区间上的反常积分}

定义6.2\quad 设对任何大于$a$的实数$b$，$f(x)$在$[a,b]$上均可积，则称极限
\[
\lim_{b\to+\infty}\int_a^bf(x)dx
\]
为$f(x)$在无穷区间$[a,+\infty)$上的反常积分（或广义积分），记为$\int_a^{+\infty}f(x)dx$，即

当这个极限存在时，则称反常积分$\int_a^{+\infty}f(x)dx$收敛（存在），否则称它发散。

类似地，定义反常积分


其中$c$为任一实常数。反常积分$\int_{-\infty}^{+\infty}f(x)dx$收敛的充要条件是两个反常积分$\int_{-\infty}^{c}f(x)dx$和$\int_{c}^{+\infty}f(x)dx$均收敛。

若$F(x)$是连续函数$f(x)$的原函数，计算反常积分时，为书写方便，记
\[
\begin{aligned}
F(+\infty) &= \lim_{x \to +\infty} F(x), \quad F(-\infty) = \lim_{x \to -\infty} F(x), \\
\int_a^{+\infty} f(x) \, dx &= F(x) \bigg|_a^{+\infty} = F(+\infty) - F(a), \\
\int_{-\infty}^b f(x) \, dx &= F(x) \bigg|_{-\infty}^b = F(b) - F(-\infty), \\
\int_{-\infty}^{+\infty} f(x) \, dx &= F(x) \bigg|_{-\infty}^{+\infty} = F(+\infty) - F(-\infty).
\end{aligned}
\]
这时反常积分的收敛与发散取决于$F(+\infty)$和$F(-\infty)$是否存在。

\subsection{无界函数的反常积分}

定义6.3\quad 若$\forall \varepsilon > 0$，$f(x)$在$[a+\varepsilon,b]$上可积，在$a$点右邻域内$f(x)$无界（称$a$为瑕点），称极限
\[
\lim_{\varepsilon\to0^+}\int_{a+\varepsilon}^bf(x)dx
\]
为无界函数$f(x)$在$(a,b]$上的反常积分（或瑕积分），记为$\int_a^bf(x)dx$，即
\[
\int_a^bf(x)dx=\lim_{\varepsilon\to0^+}\int_{a+\varepsilon}^{b}f(x)dx.
\]
当这个极限存在时，则称反常积分

\section{反常积分敛散性判别法}

反常积分敛散性判别法

比较判别法1

无穷小比较法1

比较判别法2

无穷大比较法2

$\Gamma$函数

\section{定积分的应用举例}

微元法

平面区域的面积

立体体积

平均值

平面曲线的弧长

功的计算

力与力矩的计算